% Integration notes: finite constructive Khintchine theorem.
% This file has no preamble; the root manuscript may import or adapt it.
\section{A finite quantitative Khintchine theorem}
\label{sec:finite-khintchine}

We use the normalization
\[
 \left|x-\frac pq\right|<\frac{\psi(q)}q,
 \qquad q\geq1,
\]
so that the critical series is $\sum_q\psi(q)$.  In the normalization
$|x-p/q|<\Psi(q)$, put $\psi(q)=q\Psi(q)$: the monotonicity hypothesis is
that $q\Psi(q)$ is nonincreasing and the critical series is
$\sum_q q\Psi(q)$.  All approximation functions in the finite theorem
are nonnegative rational valued.  A primitive recursive approximation
function and primitive recursive modulus give a primitive recursive
construction; no uniform evaluator for all primitive recursive programs
is presupposed.

For a finite union of intervals with rational endpoints, write $|U|$ for
its exact rational length.  It is computed by sorting endpoints and adding
lengths of the covered elementary segments.  Single endpoints have length
zero.  Every measure argument below is an identity or inequality about
these finite rational sums.

\subsection{The hypotheses and conclusions}

For $q\geq1$, let $A_q$ be the relatively open subset of $[0,1]$ consisting
of the points satisfying the displayed approximation inequality for some
integer $p$.  Numerators $0\leq p\leq q$ suffice on $[0,1]$.

\paragraph{Convergent version.}
Suppose a function $C$ satisfies $C(s)\ge1$ and
\[
 \sum_{q=C(s)}^R\psi(q)\leq2^{-s}
 \quad\hbox{for every }s\hbox{ and every }R\geq C(s).
 \tag{K-C}
\]
Then, for every $m$ and $R\geq C(m+1)$,
\[
 \left|\bigcup_{q=C(m+1)}^R A_q\right|\leq2^{-m}.
 \tag{K1}
\]
Monotonicity of $\psi$ is not needed here.

\paragraph{Divergent version.}
Suppose $\psi$ is nonincreasing and a unit-tail modulus $D$ satisfies
\[
 D(s)\geq s,\qquad
 \sum_{q=s}^{D(s)}\psi(q)\geq1 \quad(s\geq1).
 \tag{K-D}
\]
There is a primitive recursive construction, relative to the fixed
primitive recursive functions $\psi,D$, that takes $N\geq1$ and $m\geq0$
and returns a finite list $U(N,m)$ of rational open intervals and an integer
$M(N,m)$ such that
\[
 U(N,m)\subseteq\bigcup_{q=N}^{M(N,m)}A_q,
 \qquad |[0,1]\setminus U(N,m)|<2^{-m}.
 \tag{K2}
\]
Every interval in the returned list is labelled by a reduced fraction
$p/q$ with $q>N$, and is contained in its approximation interval.  Thus the
construction also proves the version requiring coprime numerators and
denominators.

A usual primitive recursive initial-sum divergence rate $H$, with
$\sum_{q=1}^{H(t)}\psi(q)\geq t$, supplies (K-D): take
\[
 D(s)=\max\left(s,H\left(1+\left\lceil
       \sum_{q=1}^{s-1}\psi(q)\right\rceil\right)\right).
\]
All quantities in this formula are rational or natural numbers, and its
verification uses only nonnegativity and subtraction of finite sums.

\subsection{Convergence: the complete estimate}

Put $r=\psi(q)/q$.  The intervals centered at $p/q$, $1\leq p\leq q-1$,
have clipped length at most $2r$; the two intervals at $0$ and $1$ have
clipped length at most $r$ each.  Therefore
\[
 |A_q|\leq2qr=2\psi(q).
\]
Finite subadditivity and (K-C) immediately give (K1).  The same calculation
bounds every finite initial portion of the countable tail cover, which is
the constructive null-cover meaning of the convergence conclusion.  If
one insists on intervals open in all of $\mathbb R$ and counts their
unclipped lengths, the simpler bound $4\psi(q)$ applies; replace
$C(m+1)$ by $C(m+2)$.

\subsection{Enough separated reduced fractions in every interval}

Let $J=(a,b)\subseteq(0,1)$, allowing $a=0$ or $b=1$, and put $L=b-a>0$.
Its middle half is
\[
 J^*=(a+L/4,b-L/4),\qquad \lambda=|J^*|=L/2.
\]
For $Q=2^k$, $k\geq2$, let $\mathcal R_k(J)$ be the reduced fractions
$p/q\in J^*$ with $Q/4<q\leq Q$.  We prove that
\[
 \frac{LQ^2}{32}\leq\#\mathcal R_k(J)\leq LQ^2
 \tag{K3}
\]
whenever
\[
 k\geq K(J,N):=
 \max\left(4,512\lceil1/L\rceil,2+\lceil\log_2N\rceil\right).
 \tag{K4}
\]
Here $\lceil\log_2N\rceil$ is only notation for the least $j\leq N$
with $2^j\geq N$; it is a bounded search.  Its presence also ensures
$q>Q/4\geq N$ for the fractions used below.

Let $A(X)$ count all integer pairs $(p,q)$ with $1\leq q\leq X$ and
$p/q\in J^*$, without requiring coprimality.  An open interval of length
$t$ contains between $t-1$ and $t+1$ integers.  Consequently
\[
 \frac{\lambda X^2}{2}-X\leq A(X)
 \leq\frac{\lambda X^2}{2}+\left(\frac\lambda2+1\right)X.
 \tag{K5}
\]
Every noncoprime pair counted by $A(Q)$ is divisible by some integer
$d\in\{2,\ldots,Q\}$.  Dividing both coordinates by $d$ and applying
(K5) bounds the number of bad pairs by
\[
 \sum_{d=2}^Q A(\lfloor Q/d\rfloor)
 \leq\frac{\lambda Q^2}{2}\sum_{d=2}^Q\frac1{d^2}
      +\left(\frac\lambda2+1\right)Q\sum_{d=2}^Q\frac1d.
\]
No prime-number theorem, totient asymptotic, or inversion formula occurs
in this estimate.  The required elementary inequalities are
\[
 \sum_{d=2}^Q\frac1{d^2}\leq\frac14+
       \sum_{d=3}^Q\frac1{d(d-1)}\leq\frac34,
 \qquad
 \sum_{d=2}^{2^k}\frac1d\leq k.
 \tag{K6}
\]
The first telescopes, and the second splits the sum into the $k$ blocks
$(2^{j-1},2^j]$, each of sum at most one.  Subtracting the bad pairs and
then all pairs with denominator at most $Q/4$ gives
\begin{align*}
 \#\mathcal R_k(J)
 &\geq \frac{3\lambda Q^2}{32}
 -Q\left(1+\frac\lambda8+\frac14+
                 \left(\frac\lambda2+1\right)k\right)\\
 &\geq\frac{3LQ^2}{64}-3Q(k+1).
 \tag{K7}
\end{align*}
For $k\geq4$, induction gives $2^k\geq k^2$.  Condition (K4) gives
$Lk\geq512$, so
\[
 LQ\geq Lk^2\geq256(k+1)\geq192(k+1).
\]
The last term in (K7) is therefore at most $LQ^2/64$, proving the lower
bound in (K3).

Distinct reduced fractions with denominators at most $Q$ are separated
by at least $Q^{-2}$: their difference has a nonzero integer numerator
over a denominator at most $Q^2$.  In an interval of length $L/2$ this
gives at most $LQ^2/2+1$ fractions.  Our bounds imply $LQ^2\geq2$, so
this is at most $LQ^2$, proving the upper bound in (K3).

\subsection{Small approximation intervals and their intersections}

For $k\geq K(J,N)$ put
\[
 t_k=\min\{2^k\psi(2^k),1\},\qquad
 r_k=\frac{t_k}{8\,4^k},
\]
and let $E_k(J)$ be the union of the intervals of radius $r_k$ centered
at the fractions in $\mathcal R_k(J)$.  An interval with radius zero is
understood as empty.  All these intervals lie in $J$, since their centers
lie in $J^*$ and $r_k\leq1/(8Q^2)\leq L/4$.  They are pairwise disjoint
at each fixed scale, by the separation just proved.  Moreover, for their
labels $p/q$,
\[
 r_k\leq\frac{\psi(Q)}{8Q}\leq\frac{\psi(q)}q,
\]
where the last inequality uses $q\leq Q$ and monotonicity.  Thus every
$E_k(J)$ consists of valid approximations with denominators $q>N$.
Using (K3) and $2r_k=t_k/(4Q^2)$ yields
\[
 \frac{Lt_k}{128}\leq |E_k(J)|\leq\frac{Lt_k}{4}.
 \tag{K8}
\]

Here is the complete overlap estimate.  Fix $k<\ell$ and one interval
$I$ of $E_k(J)$.  A radius-$r_\ell$ interval meeting $I$ has its center
in an interval of length $|I|+2r_\ell$.  Separation at scale $\ell$
bounds the number of such centers by
$(|I|+2r_\ell)4^\ell+1$.  Multiplying by $2r_\ell$ gives
\[
 |I\cap E_\ell(J)|
 \leq\frac{t_\ell}{4}|I|
       +\left(\frac{t_\ell}{4}+1\right)2r_\ell
 \leq\frac{t_\ell}{4}|I|+\frac54(2r_\ell).
\]
Sum this over the disjoint intervals of $E_k(J)$, use (K8), and use
$\#\mathcal R_k(J)\leq L4^k$.  The result is
\[
 |E_k(J)\cap E_\ell(J)|
 \leq\frac L{16}t_kt_\ell
       +\frac{5L}{16}4^{k-\ell}t_\ell.
 \tag{K9}
\]
This is the sole correlation estimate needed.

\subsection{A finite block covers a fixed fraction of every interval}

Take any finite block $k=K,\ldots,B$ of admissible scales and write
$S=\sum_{k=K}^B t_k$.  On the finite partition by all interval endpoints,
let $F=\sum_{k=K}^B\mathbf1_{E_k(J)}$.  The integral signs in what follows
are abbreviations for finite rational sums of segment lengths times
integer values.  Equations (K8)--(K9) give
\[
 \int F\geq LS/128
\]
and
\begin{align*}
 \int F^2
 &=\sum_k|E_k|+2\sum_{k<\ell}|E_k\cap E_\ell|\\
 &\leq L\left(\frac S4+\frac{S^2}{16}+\frac{5S}{24}\right)
  =L\left(\frac{S^2}{16}+\frac{11S}{24}\right).
 \tag{K10}
\end{align*}
We used only $2\sum_{k<\ell}t_kt_\ell\leq S^2$ and the finite geometric
sum $\sum_{k<\ell}4^{k-\ell}\leq1/3$.  When $S\geq1$, the right side
of (K10) is at most $LS^2$.

For completeness, finite weighted Cauchy--Schwarz says
\[
 (\int F)^2\leq |\{F>0\}|\int F^2.
\]
On the endpoint partition this is the elementary identity that the
difference between the right and left sides is
$\sum_{i<j}w_iw_j(v_i-v_j)^2\geq0$, after omitting the cells with value
zero.  Therefore, if $S\geq1$,
\[
 \left|\bigcup_{k=K}^B E_k(J)\right|\geq\frac{L}{16384}.
 \tag{K11}
\]
All divisions have explicitly positive rational denominators.

\subsection{One call to the divergence modulus supplies the block}

Let $K=K(J,N)$, let $s=2^K$, and put
\[
 B=\lceil\log_2D(s)\rceil.
\]
The dyadic blocks $[2^k,2^{k+1}-1]$, $K\leq k\leq B$, include every
integer from $s$ to $D(s)$.  Monotonicity gives
\[
 1\leq\sum_{q=s}^{D(s)}\psi(q)
 \leq\sum_{k=K}^B 2^k\psi(2^k).
\]
For any finite nonnegative rational list $(a_k)$,
\[
 \sum_k\min(a_k,1)\geq\min\left(\sum_k a_k,1\right).
\]
This follows by induction from its two-term version, or by a bounded
case distinction according as some $a_k\geq1$.  Hence $S\geq1$, as
required for (K11).  The cap at one and the factor eight in the radii
are useful: one unit of tail mass suffices, even if concentrated at one
dyadic scale.

\subsection{Amplification by the finitely many residual gaps}

Set $b_0=16384$.  Start with the empty union $U_0$.  At each stage, write
$[0,1]\setminus U_j$ as a finite list of closed rational gaps.  For each
gap of positive length, apply the block construction to its interior
$J$ and add its block union to $U_j$.  Leave singleton gaps in the
complement list.  Freeze the old gap list at the beginning of the round:
process each old gap once, aggregate the new intervals, and only then
compute the next gap list.  All added intervals are inside the gap in which they
were constructed.  If $R_j=|[0,1]\setminus U_j|$, (K11) gives
\[
 R_{j+1}\leq(1-1/b_0)R_j.
\]
The rational binomial inequality
$\bigl(b_0/(b_0-1)\bigr)^{b_0}\geq1+b_0/(b_0-1)>2$
implies $(1-1/b_0)^{b_0}<1/2$.  After
\[
 T=b_0(m+1)
\]
stages we therefore have $R_T\leq2^{-(m+1)}<2^{-m}$.  The maximum of
the finitely many denominator bounds $2^B$ is $M(N,m)$.  This proves
(K2), with a fully specified finite algorithm.

Every loop has an explicit finite bound.  At a scale $k$, enumerate
$1\leq q\leq2^k$ and $0\leq p\leq q$, test coprimality and interval
membership by rational comparisons, and retain the required fractions.
The number of stages is prescribed in advance.  The finite endpoint
sort, gap construction, block endpoint, and new denominator maximum are
primitive recursive in the finite data.  The bounds are intentionally
enormous; efficiency is not the claim.

For membership-sensitive uses, open components that merely touch must
not be merged across their common endpoint.  Merge strict overlaps,
and retain every closed residual gap, including singleton gaps.  Their
zero length may be ignored only in the length calculation.

\subsection{What ``almost every'' means constructively}

The intrinsic iPRA result is (K1) or (K2), expressed in codes for finite
rational data.  Ordinary Lebesgue measure is an external interpretation:
(K1) makes the limsup null in the convergent case, and (K2) makes every
tail union conull, hence their countable intersection conull in the
divergent case.  No countable additivity, density theorem, or zero--one
law was used to prove the finite bounds.

There is also a useful constructive null-test presentation of the
divergent conclusion.  Fix $m$, and for each $n\geq0$ construct
$U(n+1,m+n+3)$.  Its finite closed residual-gap list has total length
less than $\eta_n=2^{-(m+n+3)}$.  If the list has $r_n\geq1$ entries,
enlarge every gap $(a,b)$, including singleton gaps, to
\[
 (a-\delta_n,b+\delta_n),\qquad
 \delta_n=\eta_n/(2r_n).
\]
The total length of this finite list is less than $2\eta_n$.  Concatenate
these lists, or use a primitive recursive pairing of the two list
indices with empty entries for invalid inner indices.  This produces a
primitive recursive open cover $V_m$ whose every finite sublist has
length sum bounded by
\[
 \sum_{n\geq0}2\eta_n=2^{-(m+1)}<2^{-m};
\]
the displayed infinite sum is shorthand for its geometric finite-prefix
bound.

If a represented real avoids $V_m$ for some $m$, it has, for every $N$,
a rational approximation with denominator greater than $N$.  Indeed the
positive rational enlargement margin separates it from every residual
gap of the relevant finite union.  A sufficiently precise rational
approximation to the real therefore locates it inside that union, and
finite interval searching supplies one of the labelled approximation
intervals.  This uses only finite rational case distinctions.  For a
completely explicit extraction bound, subdivide by the finitely many
endpoints, compute the least positive gap between distinct endpoints,
and use a real-name accuracy smaller than one quarter of that gap and
one quarter of the enlargement margin; finite local cases then choose
an original interval with positive slack.  Zero-width endpoint cells
inside the union are covered by an interval crossing the endpoint;
those outside the union belong to the enlarged residual list.

One must not intuitionistically identify the negation of
$\forall N\,\exists q\geq N$ with $\exists N\,\forall q\geq N$.
The finite-coverage statement and the null-test avoidance implication
make no such inference.  Nor do we infer a primitive recursive modulus
from bare convergence or divergence: the relevant modulus is supplied
as part of the constructive input.
